\chapter{ไฟฟ้ากระแสตรงและวงจรไฟฟ้า}
\label{ch:ch12_current_electricity}
\index{ไฟฟ้ากระแสตรงและวงจร}

\begin{conceptbox}[title=วัตถุประสงค์การเรียนรู้ประจำบท]
เมื่อสิ้นสุดการศึกษาบทเรียนนี้ นักศึกษาสามารถ
\begin{enumerate}
    \item เมื่อศึกษาบทนี้จบแล้ว ผู้เรียนควรสามารถ
    \item อธิบายความหมายของกระแสไฟฟ้า ความต่างศักย์ไฟฟ้า และความสัมพันธ์ระหว่างปริมาณทั้งสองได้
    \item นิยามความต้านทานไฟฟ้าและอธิบายกฎของโอห์มได้อย่างถูกต้อง พร้อมทั้งประยุกต์ใช้ในการคำนวณ
    \item วิเคราะห์และคำนวณความต้านทานรวม กระแสไฟฟ้า และความต่างศักย์ในวงจรไฟฟ้าแบบอนุกรมและแบบขนานได้
    \item อธิบายและประยุกต์ใช้กฎของเคอร์ชอฟฟ์ทั้งกฎกระแสและกฎแรงดันในการวิเคราะห์วงจรเบื้องต้นได้
\end{enumerate}
\end{conceptbox}

\section{ทฤษฎีและหลักการพื้นฐาน}

\begin{bioappbox}
\textbf{การนำสัญญาณประสาทและเครื่องมือวัดทางชีวการแพทย์}

การส่งสัญญาณประสาท (Action Potential) ไปตามแอกซอนของเซลล์ประสาทสามารถวิเคราะห์ได้ด้วยแบบจำลองวงจรไฟฟ้า RC ขนานตามทฤษฎีสายเคเบิล (Cable Theory) นอกจากนี้ หลักการของสะพานวีตสโตน (Wheatstone Bridge) และกฎของโอห์ม $V = IR$ ยังเป็นหัวใจสำคัญของการทำงานของหัววัดความชื้นในดิน หัววัดสภาพนำไฟฟ้าของน้ำ (EC Meter) และเครื่องตรวจคลื่นไฟฟ้าหัวใจ (ECG) ในงานวิทยาศาสตร์ชีวภาพ
\end{bioappbox}

\begin{labbox}
1. ตรวจสอบตำแหน่งศูนย์ (Zero Error) ของเครื่องมือวัดทุกชนิดก่อนเริ่มบันทึกข้อมูล
2. ทำการวัดซ้ำอย่างน้อย 3 ถึง 5 ครั้งในแต่ละจุดทดลองเพื่อคำนวณหาค่าเฉลี่ยและค่าเบี่ยงเบนมาตรฐาน
3. บันทึกผลการวัดพร้อมระบุหน่วยเอสไอ (SI Units) และค่าความคลาดเคลื่อนของการวัดเสมอ
\end{labbox}

\section{ตัวอย่างโจทย์และการวิเคราะห์เชิงฟิสิกส์}

\begin{examplebox}
ตัวอย่างที่ 12.1 หลอดไฟดวงหนึ่งต่อกับแบตเตอรี่ที่มีความต่างศักย์ 12 โวลต์ หากหลอดไฟมีความต้านทาน 24 โอห์ม จงหากระแสไฟฟ้าที่ไหลผ่านหลอดไฟ

วิธีทำ จากกฎของโอห์ม V = IR จัดรูปสมการเพื่อหากระแสไฟฟ้า จะได้ I = V/R

แทนค่า I = 12/24 = 0.5 แอมแปร์

ตอบ: กระแสไฟฟ้าที่ไหลผ่านหลอดไฟมีค่าเท่ากับ 0.5 แอมแปร์
\end{examplebox}

\section{แบบฝึกหัดทบทวนและคำถามท้ายบท}

\begin{enumerate}
    \item จงอธิบายความแตกต่างระหว่างกระแสตรง (DC) และกระแสสลับ (AC)
    \item ลวดตัวนำเส้นหนึ่งมีประจุไฟฟ้า 15 คูลอมบ์ไหลผ่านในเวลา 5 วินาที จงหากระแสไฟฟ้าที่ไหลผ่านลวดตัวนำนี้
    \item จงอธิบายความหมายของกฎของโอห์ม พร้อมเขียนสมการและระบุหน่วยของแต่ละปริมาณ
    \item ตัวต้านทานตัวหนึ่งมีความต่างศักย์คร่อม 9 โวลต์ และมีกระแสไฟฟ้าไหลผ่าน 0.3 แอมแปร์ จงหาค่าความต้านทานของตัวต้านทานนี้
    \item ตัวต้านทานสี่ตัวมีค่า 5, 10, 15 และ 20 โอห์ม ต่อกันแบบอนุกรม จงหาความต้านทานรวมของวงจร
    \item ตัวต้านทานสามตัวมีค่าเท่ากันตัวละ 12 โอห์ม ต่อกันแบบขนาน จงหาความต้านทานรวมของวงจร
\end{enumerate}

% สิ้นสุดบทเรียน
